\begin{ouverture}[geometrie]{15}{Droites, cercles et constructions}{Partie 4 · Espace et géométrie}
{Segments · Perpendiculaires · Parallèles · Cercles}
\objectif{Employer le vocabulaire et les codes de la géométrie.}
\objectif{Tracer perpendiculaires et parallèles avec les instruments.}
\objectif{Décrire un cercle par la distance au centre.}
\prerequis{Règle graduée, équerre et compas.}
\end{ouverture}
\pagestyle{otMain}\essentiel\fontsize{11}{15}\selectfont
\phantomsection\label{cours-15-1}
\essSection{1}{Les mots et les outils de la géométrie}
\begin{definition}Par deux points A et B passent : la droite (AB), illimitée des deux côtés ; le segment [AB], limité par A et B ; la demi-droite [AB), qui part de A et traverse B sans s’arrêter. La longueur du segment [AB] se note AB, sans crochets.\end{definition}
\begin{savoirfaire}[Placer un milieu]Le milieu du segment [AB] est le point du segment situé à égale distance de A et de B. Pour AB $=7$ cm, le milieu est à $3{,}5$ cm de chaque extrémité : règle graduée ou compas permettent de le placer.\end{savoirfaire}
\begin{propriete}Sur une figure, des petits traits identiques codent des longueurs égales et un petit carré code un angle droit. Les codes affirment les propriétés : une figure à main levée n’est pas une preuve.\end{propriete}
\begin{attention}La règle sert à tracer et à mesurer, l’équerre contrôle les angles droits, le compas reporte des longueurs et trace les cercles. Vérifier une propriété se fait avec l’instrument adapté.\end{attention}
\clearpage
\phantomsection\label{cours-15-2}
\essSection{2}{Droites particulières et cercles}
\begin{definition}Deux droites perpendiculaires se coupent en formant un angle droit : l’équerre le vérifie. Deux droites parallèles ne se coupent jamais : leur écart reste constant, comme les deux bords d’une règle.\end{definition}
\label{fig-cours-15-2-cours-cercle}\figureGeo{cours-cercle}

\begin{definition}Le cercle de centre O et de rayon 3 cm est formé de tous les points situés à exactement 3 cm de O. Le disque contient aussi les points intérieurs, à moins de 3 cm. Un diamètre traverse le cercle en passant par le centre : il mesure deux rayons.\end{definition}
\begin{savoirfaire}[Tracer avec l’équerre]Pour tracer la perpendiculaire à une droite passant par un point : poser un côté de l’angle droit de l’équerre sur la droite, faire glisser jusqu’au point, tracer. Pour une parallèle : équerre tenue sur la droite, règle posée contre l’équerre, faire glisser l’équerre le long de la règle.\end{savoirfaire}

\clearpage\rubrique{Découverte}
\begin{blocLivre}{ex-15-1}\begin{exercice}[1]{Les bons noms}
Sur la figure, nommer avec la bonne notation : le segment qui joint A et B ; la droite qui passe par C et D ; la demi-droite qui part de E et passe par F.
\label{fig-ex-15-1}\figureGeo{geo-1}
\end{exercice}\end{blocLivre}
\begin{blocLivre}{ex-15-2}\begin{exercice}[1]{Le milieu du segment}
Le segment [AB] mesure 7 cm. Placer son milieu M et donner la longueur AM.
\label{fig-ex-15-2}\figureGeo{geo-2}
\end{exercice}\end{blocLivre}
\begin{blocLivre}{ex-15-3}\begin{exercice}[1]{Perpendiculaires ?}
Avec l’équerre, dire quelles paires de droites sont perpendiculaires.
\label{fig-ex-15-3}\figureGeo{geo-3}
\end{exercice}\end{blocLivre}
\begin{blocLivre}{ex-15-4}\begin{exercice}[1]{Parallèles ?}
Dire si les droites tracées sont parallèles, en contrôlant l’écart avec la règle graduée.
\label{fig-ex-15-4}\figureGeo{geo-4}
\end{exercice}\end{blocLivre}
\begin{blocLivre}{ex-15-5}\begin{exercice}[1]{Le vocabulaire du cercle}
Sur le cercle de centre O, nommer un rayon, un diamètre et une corde tracés sur la figure.
\label{fig-ex-15-5}\figureGeo{geo-5}
\end{exercice}\end{blocLivre}

\rubrique{Entraînement}
\begin{blocLivre}{ex-15-6}\begin{exercice}[2]{Tracer une perpendiculaire}
Tracer la droite perpendiculaire à la droite (d) passant par le point A, qui n’est pas sur (d).
\label{fig-ex-15-6}\figureGeo{geo-6}
\end{exercice}\end{blocLivre}
\begin{blocLivre}{ex-15-7}\begin{exercice}[2]{Tracer une parallèle}
Tracer la droite parallèle à la droite (d) passant par le point B.
\label{fig-ex-15-7}\figureGeo{geo-7}
\end{exercice}\end{blocLivre}
\begin{blocLivre}{ex-15-8}\begin{exercice}[2]{Rayon et diamètre}
Un cercle a un rayon de 4 cm. Quelle est la longueur d’un diamètre ? Et si le diamètre mesure 7 cm, que vaut le rayon ?
\label{fig-ex-15-8}\figureGeo{geo-8}
\end{exercice}\end{blocLivre}
\begin{blocLivre}{ex-15-9}\begin{exercice}[2]{Lire les codes}
La figure porte des codes : petits traits sur [AB] et [AC], carré en A. Que peut-on affirmer sans mesurer ?
\label{fig-ex-15-9}\figureGeo{geo-9}
\end{exercice}\end{blocLivre}
\begin{blocLivre}{ex-15-10}\begin{exercice}[2]{À 3 cm du puits}
Sur le plan, O marque un puits. Colorier l’ensemble des points situés à exactement 3 cm de O, puis dire où sont les points à moins de 3 cm.
\label{fig-ex-15-10}\figureGeo{geo-10}
\end{exercice}\end{blocLivre}

\rubrique{Problèmes}
\begin{blocLivre}{ex-15-11}\begin{exercice}[3]{Un programme à exécuter}
Exécuter le programme : tracer un rectangle ABCD tel que AB $=5$ cm et BC $=3$ cm ; placer le milieu M de [AB] ; tracer le cercle de centre A passant par M.
\label{fig-ex-15-11}\figureGeo{geo-11}
\end{exercice}\end{blocLivre}
\begin{blocLivre}{ex-15-12}\begin{exercice}[3]{Deux cercles croisés}
Tracer le segment [AB] avec AB $=4$ cm, puis les cercles de centres A et B, tous deux de rayon 3 cm. Que peut-on dire des points où les cercles se coupent ?
\label{fig-ex-15-12}\figureGeo{geo-12}
\end{exercice}\end{blocLivre}
\begin{blocLivre}{ex-15-13}\begin{exercice}[3]{Reproduire sur quadrillage}
Reproduire la figure sur le quadrillage : un rectangle et une de ses diagonales.
\label{fig-ex-15-13}\figureGeo{geo-13}
\end{exercice}\end{blocLivre}
\begin{blocLivre}{ex-15-14}\begin{exercice}[3]{Perpendiculaires en chaîne}
Tracer une droite $(d_2)$ perpendiculaire à $(d_1)$, puis une droite $(d_3)$ perpendiculaire à $(d_2)$. Que remarque-t-on pour $(d_1)$ et $(d_3)$ ? Vérifier avec les instruments.
\label{fig-ex-15-14}\figureGeo{geo-14}
\end{exercice}\end{blocLivre}
\begin{blocLivre}{ex-15-15}\begin{exercice}[3]{Le plan du quartier}
Sur le plan, la rue des Lilas et la rue des Érables sont tracées, ainsi que l’avenue du Parc.\begin{questions}
\item Citer deux rues parallèles.
\item Citer deux rues perpendiculaires.
\item La fontaine est à égale distance des points F et G marqués sur l’avenue : que peut-on tracer pour la localiser ?
\end{questions}
\label{fig-ex-15-15}\figureGeo{geo-15}
\end{exercice}\end{blocLivre}

\clearpage
